\section{Integral homology and explicit monodromy}\label{sec:monodromy}

Fix the complex elliptic curve $E$ and the constant-$E$ family from
Section~\ref{sec:geometry}.  Let $n\geq3$ be odd and let $m>1$ divide
$n$.  We write
\[
 \Gamma_m=\mu_m,\qquad u:(Z,Y)\longmapsto(Z,e^{2\pi i/m}Y),
\]
so that $u$ is the chosen generator of the deck group.
Throughout this section $\Z[\Gamma_m]$ denotes the \emph{group ring}, not
the cyclotomic integer ring.  We first construct an integral basis
permuted by $u$, and then compute two monodromy operators in that same
basis.  Their reduction on character spaces will supply the residual
irreducibility used later.

\subsection{A basis obtained by cutting the torus}

Choose an oriented period basis $a,b$ and identify the underlying
oriented torus with $(\mathbb R/\Z)^2$, with $a$ horizontal and $b$ vertical.
Consider the path of parameters
\begin{equation}\label{mon:horizontal-path}
 R_s=\frac{s}{n}a,\qquad 0<s<\frac12.
\end{equation}
Its branch points are $-sa$ and $sa$.  Let $I_s$ be the horizontal arc
from $-sa$ to $sa$ through $O$.

\begin{lemma}\label{mon:slit}
For the parameters in \eqref{mon:horizontal-path}, the cover
$C_{m,R_s}\to E$ is trivial over $E\setminus I_s$.  It is obtained by
cutting $m$ copies of $E$ along $I_s$ and joining adjacent banks
cyclically.  Crossing the slit changes the sheet index by $1$ or $-1$,
according to the orientation of crossing and the initial choice of
indexing.
\end{lemma}

\begin{proof}
For small $s$, enclose the colliding sections in a small disk about
$O$.  The model at $R=O$ restricts outside that disk to $m$ disjoint
copies of its complement, by Proposition~\ref{geo:compact-type}.
If $D$ is such a disk, the inclusion
$E\setminus\operatorname{int}(D)\hookrightarrow E\setminus I_s$
is a homotopy equivalence.  The unramified cover of $E\setminus I_s$
is therefore trivial.  As $s$ varies in $(0,1/2)$, the arcs $I_s$ move by isotopy,
and the topological isomorphism class of this unramified cover remains
constant.  The monodromy around either branch point is a generator
of $\Gamma_m$, since its valuation in $f_R$ is $1$ or $-1$.
This gives the stated gluing.
\end{proof}

On each sheet let $a_i$ be the horizontal loop at vertical coordinate
$1/2$ and let $b_i$ be the vertical loop at horizontal coordinate
$1/2$, for $i\in\Z/m\Z$.  These loops avoid the slit, intersect at
one point on their common sheet, and may be oriented so that
\begin{equation}\label{mon:intersection}
 \langle a_i,b_j\rangle=\delta_{ij},\qquad
 \langle a_i,a_j\rangle=\langle b_i,b_j\rangle=0.
\end{equation}
Index the sheets so that $u$ sends $a_i,b_i$ to $a_{i+1},b_{i+1}$.

\begin{proposition}\label{mon:integral-module}
The cycles $a_i,b_i$ form an integral symplectic basis of
$H_1(C_{m,R_s},\Z)$.  Consequently, for every smooth fiber of the
family,
\begin{equation}\label{mon:free-lattices}
 H_1(C_m,\Z)\simeq\Z[\Gamma_m]^2,
 \qquad H^1(C_m,\Z)\simeq\Z[\Gamma_m]^2.
\end{equation}
In particular, after comparison, $H^1_{\mathrm{et}}(C_m,\Z_\ell)$ is
free of rank two over $\Z_\ell[\Gamma_m]$ for every prime $\ell$, including
$\ell\mid m$.
\end{proposition}

\begin{proof}
The genus is $m$, so $H_1$ has rank $2m$.  Equation~\eqref{mon:intersection}
shows that the $2m$ displayed classes are linearly independent and
span a sublattice with unimodular intersection pairing.  If its index
in $H_1$ were $d$, its intersection determinant would be $d^2$ times
the determinant on $H_1$.  Both determinants are $1$, so $d=1$.
The deck action now gives the first isomorphism in
\eqref{mon:free-lattices}.

The dual of the regular integral module is again regular: explicitly,
the coefficient-of-identity pairing
\[
 \Z[\Gamma_m]\times\Z[\Gamma_m]\longrightarrow\Z,
 \qquad (x,y)\longmapsto [1](xy)
\]
is perfect.  With the pullback action on the dual, it identifies
$\operatorname{Hom}_{\Z}(\Z[\Gamma_m],\Z)$ with $\Z[\Gamma_m]$.
This proves the cohomological assertion without dividing by $m$.
Topological local triviality transports the result through the connected
smooth parameter open, and comparison gives the \'etale assertion.
\end{proof}

\begin{remark}
The compact-type degeneration gives another explanation of the integral
statement.  Its Jacobian extends to an abelian scheme with special
fiber $E^m$, and $\Gamma_m$ permutes the factors regularly.
The integral homology of this fiber is the regular module twice.
The slit construction makes this description concrete by identifying
cycles that can also be followed towards the other degeneration.
\end{remark}

\subsection{Vanishing cycles and a common monodromy basis}

Let $s$ approach $1/2$ in \eqref{mon:horizontal-path}.  The branch
points meet at $a/2$, while the slit fills all but a shrinking gap in
the horizontal circle.  Choose a small disk about $a/2$ containing
both branch points, with radius larger than their distance from the
collision point.  Its boundary has $m$ closed lifts $\delta_i$.
They are the neck circles at the $m$ nodes of
Lemma~\ref{geo:collision}; hence they are the vanishing cycles.

\begin{lemma}\label{mon:vanishing-classes}
After a cyclic indexing of the lifts, their homology classes are
\begin{equation}\label{mon:delta}
 [\delta_i]=\pm(a_i-a_{i+1}),\qquad i\in\Z/m\Z,
\end{equation}
or the same formula with $i-1$ in place of $i+1$.
\end{lemma}

\begin{proof}
Arrange the boundary of the disk to cross the slit transversely twice.
A lifted boundary circle has its upper arc on one sheet and its lower
arc on an adjacent sheet.  Index the upper sheet by $i$; the lower
one is $i+1$ or $i-1$, by Lemma~\ref{mon:slit}.  The circle is disjoint
from every $a_j$.  It crosses the two corresponding vertical loops
once each, with opposite intersection signs, and no other $b_j$.
These intersections determine its class in the integral symplectic
basis of Proposition~\ref{mon:integral-module}, giving
\eqref{mon:delta}.  See Figure~\ref{mon:slit-figure}.
\end{proof}

\begin{figure}[tb]
\centering
\begin{tikzpicture}[
  x=1cm,y=1cm,
  every node/.style={font=\small,inner sep=1.5pt},
  boundary/.style={draw=black!45,line width=.5pt},
  slit/.style={draw=black!75,line width=1.7pt},
  cycle/.style={draw=blue!60!black,line width=.8pt},
  arrow/.style={-{Stealth[length=1.5mm,width=1.1mm]}}
]
  \node at (2,4.65) {\textup{(a)} Slit in the base torus};
  \node at (9,4.65) {\textup{(b)} A lifted boundary circle};

  \draw[boundary] (0,0) rectangle (4,4);
  \draw[boundary,arrow] (.55,0) -- (1,0);
  \draw[boundary,arrow] (.55,4) -- (1,4);
  \draw[boundary,arrow] (0,2) -- (0,2.4);
  \draw[boundary,arrow] (4,2) -- (4,2.4);
  \draw[boundary,arrow] (0,2.1) -- (0,2.5);
  \draw[boundary,arrow] (4,2.1) -- (4,2.5);
  \node[below] at (0,0) {$0$};
  \node[below] at (4,0) {$1$};
  \node[below] at (2,-.32) {horizontal coordinate $x$};

  \draw[cycle] (0,3) -- (4,3);
  \draw[cycle,arrow] (.8,3) -- (1.2,3);
  \node[above] at (3.35,3) {$a$};
  \draw[cycle] (2,0) -- (2,4);
  \draw[cycle,arrow] (2,3.45) -- (2,3.8);
  \node[right] at (2,3.65) {$b$};

  \draw[slit] (0,1) -- (1.52,1);
  \draw[slit] (2.48,1) -- (4,1);
  \fill (1.52,1) circle (1.7pt);
  \fill (2.48,1) circle (1.7pt);
  \node[left] at (0,1) {$O$};
  \node[above] at (.65,1) {slit};
  \draw[densely dashed,black!60,line width=.55pt] (2,1) circle (.72);
  \node[below left,fill=white] at (1.52,1) {$+s$};
  \node[below right,fill=white] at (2.48,1) {$-s$};
  \node[above right,font=\footnotesize] at (2.6,1.58)
    {$\partial D$};
  \node[right,font=\footnotesize] at (2,1.15) {$a/2$};

  \draw[slit] (6.9,2) -- (8.58,2);
  \draw[slit] (9.42,2) -- (11.1,2);
  \fill (8.58,2) circle (1.7pt);
  \fill (9.42,2) circle (1.7pt);
  \node[below] at (8.58,2) {$+s$};
  \node[below] at (9.42,2) {$-s$};
  \draw[cycle] (7.55,2) arc[start angle=180,end angle=0,radius=1.45];
  \draw[cycle,dashed] (10.45,2) arc[start angle=0,end angle=-180,radius=1.45];
  \draw[cycle,arrow] (8.1683,3.1878)
    arc[start angle=125,end angle=105,radius=1.45];
  \draw[cycle,arrow] (10.1107,1.0679)
    arc[start angle=-40,end angle=-60,radius=1.45];
  \node[above left] at (8.15,3.2) {sheet $i$};
  \node[below left] at (8.15,.82) {sheet $i+1$};
  \node[right] at (10.45,2.7) {$\delta_i$};

  \draw[black!70,line width=.65pt,arrow] (9,2.78) -- (9,4.03);
  \draw[black!70,line width=.65pt,arrow] (9,-.05) -- (9,1.22);
  \node[right] at (9,3.9) {$b_i$};
  \node[right] at (9,.12) {$b_{i+1}$};
  \node[left,fill=white,inner sep=1pt] at (8.96,3.48) {$+$};
  \node[left,fill=white,inner sep=1pt] at (8.96,.5) {$-$};
  \node at (9,-.65) {$[\delta_i]=a_i-a_{i+1}$};
\end{tikzpicture}

\caption{The horizontal slit leaves a gap about the collision point
$a/2$.  A small circle about that gap crosses the slit twice.  Its
lift has one arc on sheet $i$ and the other on sheet $i+1$; its two
intersections with the lifted vertical cycles have opposite signs.
This identifies the vanishing class as $\pm(a_i-a_{i+1})$.
Opposite edges of the torus rectangle are identified; its coordinates
range over $[0,1]\times[-1/4,3/4]$, so $-s$ is plotted at $1-s$.
In the right panel, solid and dashed arcs distinguish the two sheets.}
\label{mon:slit-figure}
\end{figure}

\begin{theorem}\label{mon:transvections}
In one integral basis $a_i,b_i$ as above, the geometric monodromy of
the constant-$E$ family contains operators $T_a,T_b$ satisfying
\begin{align}
 T_a(a_j)&=a_j,&
 T_a(b_j)&=b_j+\sigma_a(2a_j-a_{j-1}-a_{j+1}),
 \label{mon:Ta}\\
 T_b(b_j)&=b_j,&
 T_b(a_j)&=a_j+\sigma_b(2b_j-b_{j-1}-b_{j+1}),
 \label{mon:Tb}
\end{align}
where $\sigma_a,\sigma_b\in\{1,-1\}$ are independent of $j$.
Equivalently, their two off-diagonal entries, in the basis of the free
$\Z[\Gamma_m]$-module, are $\sigma_a(2-u-u^{-1})$ and
$\sigma_b(2-u-u^{-1})$.
The geometric monodromy of the family over $U$ contains these actions
as well, under the pullback identification.
\end{theorem}

\begin{proof}
The Picard--Lefschetz formula for a positively oriented loop about a
semistable fiber is the product of the Dehn twists at its nodes;
see~\cite[Expos\'e~XV]{SGA7II}.  In the present family, every node
has smoothing order one by \eqref{geo:node-equation}, so every twist
has exponent one.  The vanishing cycles are disjoint.  Their
intersection pairings with one another therefore vanish, and the
product acts on a class $\gamma$ by
\begin{equation}\label{mon:PL}
 \gamma\longmapsto
 \gamma+\sigma\sum_i
 \langle\gamma,\delta_i\rangle[\delta_i],
\end{equation}
with one sign $\sigma$ depending only on the intersection convention.
Substituting \eqref{mon:delta} gives zero on the $a_j$ and, up to this
one sign,
\[
 (a_j-a_{j+1})-(a_{j-1}-a_j)
   =2a_j-a_{j-1}-a_{j+1}
\]
on $b_j$.  This proves \eqref{mon:Ta}.  Replacing the horizontal
path by $R_s=(s/n)b$ gives \eqref{mon:Tb}, initially in the basis
obtained from the corresponding vertical slit.

It remains to compare the two bases; this matters because separate
triangular matrices in unrelated bases would not give the stated
conclusion.  Choose the starting parameters of both paths sufficiently
close to $O$.  The horizontal loop at $y=1/2$ and the vertical loop
at $x=1/2$ avoid one fixed small disk about $O$ and avoid both short
slits.  Outside this disk the covering extends across $R=O$ as $m$
disjoint sheets.  Trivialize it over a small, simply connected
parameter disk containing $O$, and use this trivialization to label
the lifts of the two fixed loops.  Parallel transport between the
starting parameters now identifies both sets of $a_i,b_i$ with the
same basis.  There is no monodromy ambiguity from circling $O$ on
these classes: the sheet trivialization extends over $O$ itself.
A different initial labeling makes only a common cyclic shift, which
leaves \eqref{mon:Ta}--\eqref{mon:Tb} unchanged.

Both loops therefore define the asserted operators in a single
geometric monodromy group.  Pullback along the open parameter curve
in $E\to U$ maps its fundamental group into that of $U$, so these
operators occur in the latter image too.  Removing additional finitely
many points from either smooth parameter open induces a surjection
on fundamental groups and does not remove any of the operators.
\end{proof}

The operator $2-u-u^{-1}$ is the cyclic difference operator arising
from the adjacent-sheet vanishing classes.  The next sections use
the integral module in Proposition~\ref{mon:integral-module} to
specialize rotation characters, and the two matrices above to keep
the resulting residual representations irreducible.
